% appP_register.tex -- generated by writing/tools/blueprint_gen.py 1.0 (level 0) on 2026-10-14 from sources 08ca495089; do not edit
\providecommand{\bpNum}[1]{\hyperref[#1]{\ref*{#1}}}%
\providecommand{\bpNumA}[1]{\hyperref[#1]{\textcolor{orange!65!black}{\ref*{#1}}}}%
\providecommand{\bpPage}[1]{\pageref*{#1}}%
\providecommand{\bpLab}[1]{{\ttfamily #1}}%
\providecommand{\bpBrk}{\allowbreak}%
\providecommand{\bpT}[2]{\facttagbox[#1]{\scriptsize #2}}%
\providecommand{\bpU}[1]{{\color{gray!70!black}#1}}%

% heading omitted: sec:appP-register is defined in a unit file

本表由各单元的 \texttt{assumption} 环境与假设登记表自动生成，按所属单元的书序排列；“被使用于”一列由各陈述的 \texttt{\string\uses} 数据计算（直接使用者的编号），与环境中「被使用于」字段的比较见检查 C-P10。
\begingroup\footnotesize\setlength{\tabcolsep}{2.5pt}\renewcommand{\arraystretch}{1.06}
\begin{longtable}{@{}>{\raggedright\arraybackslash}p{\dimexpr(\textwidth-12\tabcolsep-0mm)*149/1000\relax}>{\raggedright\arraybackslash}p{\dimexpr(\textwidth-12\tabcolsep-0mm)*84/1000\relax}>{\raggedright\arraybackslash}p{\dimexpr(\textwidth-12\tabcolsep-0mm)*153/1000\relax}>{\raggedright\arraybackslash}p{\dimexpr(\textwidth-12\tabcolsep-0mm)*127/1000\relax}>{\raggedright\arraybackslash}p{\dimexpr(\textwidth-12\tabcolsep-0mm)*133/1000\relax}>{\raggedright\arraybackslash}p{\dimexpr(\textwidth-12\tabcolsep-0mm)*234/1000\relax}>{\raggedright\arraybackslash}p{\dimexpr(\textwidth-12\tabcolsep-0mm)*117/1000\relax}@{}}
\caption{编号假设登记表（由 \texttt{assumption} 环境生成）}\label{tab:appP-register}\\
\toprule
\textbf{编号／标签} & \textbf{单元} & \textbf{内容} & \textbf{标签} & \textbf{地位} & \textbf{被使用于} & \textbf{解除途径} \\
\midrule
\endfirsthead
\caption[]{编号假设登记表（由 \texttt{assumption} 环境生成）（续）}\\
\toprule
\textbf{编号／标签} & \textbf{单元} & \textbf{内容} & \textbf{标签} & \textbf{地位} & \textbf{被使用于} & \textbf{解除途径} \\
\midrule
\endhead
\midrule
\multicolumn{7}{r@{}}{\footnotesize（续下页）}\\
\endfoot
\bottomrule
\endlastfoot
\bpNumA{asm:intro-wightman}\newline \bpLab{\detokenize{asm:}\bpBrk \detokenize{intro-}\bpBrk \detokenize{wightman}} & \mbox{第~1~章} & QCD 的 Wightman 结构 (W1)(W4) & \tagC & 全局非微扰假设。公开问题。 & \bpNum{thm:dis-moments} & 构造 QCD。 \\
\bpNumA{asm:intro-np-ope}\newline \bpLab{\detokenize{asm:}\bpBrk \detokenize{intro-}\bpBrk \detokenize{np-}\bpBrk \detokenize{ope}} & \mbox{第~1~章} & 非微扰算符乘积展开 (O1)(O2) & \tagC & 全局非微扰假设。公开问题。 & \bpNum{thm:dis-moments} & 非微扰定义。 \\
\bpNumA{asm:intro-bridge}\newline \bpLab{\detokenize{asm:}\bpBrk \detokenize{intro-}\bpBrk \detokenize{bridge}} & \mbox{第~1~章} & 桥接假设 & \tagC & 全局非微扰假设。公开问题。 & \bpNum{thm:dy-factorization}，\bpBrk \bpNum{prop:dy-npline} & 无。 \\
\bpNumA{asm:appN-collins}\newline \bpLab{\detokenize{asm:}\bpBrk \detokenize{appN-}\bpBrk \detokenize{collins}} & \mbox{附录~N} & Collins 条件 (CS1)(CS3) & \tagA & 全阶论证覆盖。 & \bpNum{thm:method-template}，\bpBrk \bpNum{thm:dy-factorization}，\bpBrk \bpNum{prop:dy-npline}，\bpBrk \bpNum{thm:appN-subtraction}，\bpBrk \bpNum{thm:appN-noprem}，\bpBrk \bpNum{prop:appN-late} & 一般图的证明。 \\
\bpNumA{asm:appN-values}\newline \bpLab{\detokenize{asm:}\bpBrk \detokenize{appN-}\bpBrk \detokenize{values}} & \mbox{附录~N} & 取值规则 (V1)(V2) & 无 & 约定。 & \bpNum{thm:method-template}，\bpBrk \bpNum{thm:appN-subtraction} & 无（约定）。 \\
\bpNumA{asm:appN-brst-nonlocal}\newline \bpLab{\detokenize{asm:}\bpBrk \detokenize{appN-}\bpBrk \detokenize{brst-}\bpBrk \detokenize{nonlocal}} & \mbox{附录~N} & BRST 分类 & \tagC & 无全阶论证（具名缺口）。 & \bpNum{prop:dy-ctech}，\bpBrk \bpNum{thm:appN-structure} & 分类定理。 \\
\bpNumA{asm:appN-unused}\newline \bpLab{\detokenize{asm:}\bpBrk \detokenize{appN-}\bpBrk \detokenize{unused}} & \mbox{附录~N} & 未被使用的假设 & \tagA & 全阶论证覆盖。 & — & 无。 \\
\bpNumA{asm:appN-wrongtag}\newline \bpLab{\detokenize{asm:}\bpBrk \detokenize{appN-}\bpBrk \detokenize{wrongtag}} & \mbox{附录~N} & 标签不符的假设 & \tagA & 仅在特例中证明。 & \bpNum{thm:appN-structure} & 见该假设的「解除途径」 \\
\bpLab{\detokenize{asm:}\bpBrk \detokenize{appN-}\bpBrk \detokenize{missing}} & \mbox{附录~N} & 未写出的假设 & \tagA & 文献有全阶论证，未作为精确陈述证明 & — & （单元尚未写入） \\
\end{longtable}
\endgroup
